\section{Grammatik}

Die Autoren definieren in ihrem Papier eine Grammatik, die wir im Folgenden in eine EBNF-nahe Hilfsdarstellung überführen.\footnote{
    Für eine bessere Lesbarkeit verzichten wir auf eine Anwendung von Anführungszeichen für das Terminalalphabet $\Sigma$.
    Unterschiede zwischen Metasymbolen und Grammatikelementen sind außerdem farblich hervorgehoben.
}
Bei der Notation orientieren wir uns hierbei an der Definition kontextfreier Grammatiken von Vossen und Witt~\cite{VW16}:

\[
 G_{\text{CoreECS}} = (\Sigma, N, P, S)
\]

\noindent
Typbindungen und Ausdrücke der darunterliegenden Sprache fleißen auch bei uns in Anlehnung an Redmond et al. in die Definition ein.
Auf diese gehen wir nachfolgend gesondert ein.


\[
    \begin{alignedat}{3}
\Sigma &\coloneqq{}&
\{&
\\[-1mm]
&&&
\tm{a}, \tm{b}, \ldots, \tm{z},\\
&&&
\tm{incl}, \tm{excl}, \tm{anyway},\\
&&&
\tm{attach}, \tm{detach}, \tm{\nu}, \tm{nil},\\
&&&
\tm{conc}, \tm{seq},\\
&&&
\tm{(}, \tm{)}, \tm{,}, \tm{||},\\
&&&
\tm{\land}, \tm{\bullet}, \tm{\fatsemi}
\\[-1mm]
&&&
\}
\\[2mm]
N &\coloneqq{}&
\{&
\\[-1mm]
&&&
\nt{char}, \nt{label},\\
&&&
\nt{query\_command}, \nt{constraint}, \nt{query},\\
&&&
\nt{mutation\_command}, \nt{mutation},\\
&&&
\nt{system},\\
&&&
\nt{schedule\_command}, \nt{schedule}
\\[-1mm]
&&&
\}
\\[2mm]
P &\coloneqq{}&
\{&
\\[-1mm]
&&&
\begin{array}{@{}lcl@{}}
    \nt{char}
    &\Coloneqq&
    \tm{a} \mid \tm{b} \mid \ldots \mid \tm{z}
    \\[1mm]

    \nt{label}
    &\Coloneqq&
    \nt{char}\ \{\nt{char}\}
    \\[1mm]

    \nt{query\_command}
    &\Coloneqq&
    \tm{incl} \mid \tm{excl} \mid \tm{anyway}
    \\[1mm]

    \nt{constraint}
    &\Coloneqq&
    \nt{query\_command}\ \nt{label}
    \\[1mm]

    \nt{query}
    &\Coloneqq&
    \nt{constraint}\ \{\tm{\land}\ \nt{constraint}\}
    \\[1mm]


    \nt{mutation\_command}
    &\Coloneqq&
    \tm{attach}\ \nt{label}\ \tm{(}\ \ty{e}\ \tm{,}\ \ty{k_i}\ \tm{)}
    \\
    &\mid&
    \tm{detach}\ \nt{label}\ \tm{(}\ \ty{e}\ \tm{)}
    \\
    &\mid&
    \tm{\nu}\ \tm{(}\ \metaexpr{\ty{E}\ \ty{\rightarrow}\ \ty{M}}\ \tm{)}
    \\
    &\mid&
    \tm{nil}
    \\[1mm]

    \nt{mutation}
    &\Coloneqq&
    \nt{mutation\_command}\ \{\tm{\bullet}\ \nt{mutation\_command}\}
    \\[1mm]

    \nt{system}
    &\Coloneqq&
    \tm{(}\ \nt{query}\ \tm{,}\ \metaexpr{\ty{E}^{\ty{\mathrm{dim}}\ty{(}\nt{query}\ty{)}}
    \ \ty{\times}\ \ty{(}\ty{|}\nt{query}\ty{|}\ty{)}
    \ \ty{\rightarrow}\ \ty{M}}\ \tm{)}
    \\[1mm]

    \nt{schedule\_command}
    &\Coloneqq&
    \tm{conc}\ \tm{(}\ \nt{system}\ \tm{)}
    \\
    &\mid&
    \tm{seq}\ \tm{(}\ \nt{system}\ \tm{)}
    \\[1mm]

    \nt{schedule}
    &\Coloneqq&
    \nt{schedule\_command}\
    \\
    &\mid&
    \nt{schedule}\ \{\tm{||}\ \nt{schedule}\}
    \\
    &\mid&
    \nt{schedule}\ \{\tm{\fatsemi}\ \nt{schedule}\}

\end{array}
\\[-1mm]
&&&
\}
\\[2mm]
S &\coloneqq{}& \nt{schedule}
    \end{alignedat}
\]

\bigskip
\noindent
Es gelten in Anlehnung an Remdond et al. darüber hinaus weitere folgende Typeinschränkungen:

\begin{itemize}
    \item \nt{query :}\ \ty{Q} (\nt{query} erzeugt Terme vom Typ \ty{Q})
    \item \nt{mutation :}\ \ty{M} (\nt{mutation} erzeugt Terme vom Typ \ty{M})
    \item \nt{system :}\ \ty{S} (\nt{system} erzeugt Terme vom Typ \ty{S})
    \item \nt{schedule :}\ \ty{Z} (\nt{schedule} erzeugt Terme vom Typ \ty{Z})
\end{itemize}

\noindent
Des Weiteren nutzen wir Symbole, die nicht Bestandteile der Grammatik sind.
Hier greift eine von Redmond et al. implizite, nicht formalisierte Anforderung, dass diese konkreten Typen beziehungsweise Terme von der darunterliegenden Programmiersprache zur Verfügung gestellt werden sollen.

\begin{itemize}
    \item \ty{E}: Entspricht einem konkreten Entitätstyp $E$.
    \item \ty{e}: Entspricht einem Entitätsterm $e$ vom Typ \ty{E}.
    \item \ty{k_i}: Entspricht einem Komponententerm vom Typ $K_i$.
\end{itemize}

\noindent
Eine semantische Bedingung im Zusammenhang mit der Verwendung von \ty{k_i} und \nt{label} ist, dass für jede Produktion, die diese Symbole enthält, dieses Label als auch \ty{k_i} übereinstimmen müssen: Als gedankliches Modell kann man sich hier vorstellen, dass diese Produktionen an den Komponententyp - vertreten durch das Label beziehungsweise \ty{k_i} in der Grammatik - gebunden sind.\\

\noindent
Weiterhin sind durch einen andersfarbigen Hintergrund hervorgehobene Ausdrücke wie beispielsweise

\[
    \metaexpr{\ty{E}\ \ty{\rightarrow}\ \ty{M}}
\]

\noindent
wiederum Teil der darunterliegenden Programmiersprache und drücken aus, dass dies typgebundene Terme sind.
In diesem Fall wird von der darunterliegenden Programmiersprache eine Funktion erwartet, die eine Eingabe vom Typ \ty{E} entgegennimmt und einen Wert vom Typ \ty{M} zurückliefert.\\


\bigskip
\noindent
Damit schließen wir diesen Abschnitt mit der Angabe der von Redmond et {al.} angegebenen Grammatik, auf die wir uns im weiteren Verlauf größtenteils beziehen werden.
Einige Formatierungen haben wir aus unserer Hilfsdefinition zur einfacheren Gegenüberstellung wiederverwendet.

\[
    \begin{array}{@{}rcll@{}}
        q : \ty{Q}
        &\Coloneqq&
        \tm{incl}_{i}
        \mid
        \tm{excl}_{i}
        \mid
        \tm{anyway}_{i}
        \mid
        (q \land q')
        &
        \text{Query}
        \\[1mm]

        m : \ty{M}
        &\Coloneqq&
        \tm{attach}_{i}(\ty{e},\ty{k_i})
        \mid
        \tm{detach}_{i}(\ty{e})
        \mid
        (m \tm{\bullet} m')
        \mid
        \nu\!\left(\metaexpr{\ty{E} \rightarrow \ty{M}}\right)
        \mid
        \tm{nil}
        &
        \text{Mutation}
        \\[1mm]

        s : \ty{S}
        &\Coloneqq&
        \left(
            \vec{q},
            \metaexpr{
                \ty{E}^{\ty{\mathrm{dim}(}\vec{q}\ty{)}}
                \ty{\times}
                \ty{(|}\vec{q}\ty{|)}
                \ty{\rightarrow}
                \ty{M}
            }
        \right)
        &
        \text{System}
        \\[1mm]

        z : \ty{Z}
        &\Coloneqq&
        \tm{conc}(s)
        \mid
        \tm{seq}(s)
        \mid
        (z \tm{||} z')
        \mid
        (z \tm{\fatsemi} z')
        &
        \text{Schedule}
    \end{array}
\]

\noindent
Mit den Details beschäftigen sich die Abschnitte, die sich den jeweiligen Produktionen widmen.

